\section{Discussion and Remaining Evaluation}
\label{sec:discussion}
\subsection{Where the Benefit Comes From}
The primary observed benefit is eliminating repeated construction and movement of a shared formula. Hardware argmax is one way to produce the resulting small update, but the optimized CPU path also produces it cheaply. Once both paths reduce interface work to approximately a millisecond per puzzle, Linux-managed NVDLA execution dominates the approximately 2.5-second pipeline. Increasing conversion throughput alone has little remaining end-to-end opportunity in this setting.

The optimized system is therefore not yet a balanced producer--consumer pipeline. Each given is a separate NVDLA inference submitted through Linux, at approximately 79\,ms per cell, or about 2.49\,s for 31.4 givens; the producer is roughly 250--800 times slower than a 3--10\,ms solve and accounts for about 99\% of a delta run. LeNet on \texttt{nv\_small} should require far less time, so most of this interval is likely spent in NVDLA's user-mode runtime and kernel driver on the 78\,MHz Ariane core, although we have not separated device time from software time. The CPU also mediates every cell: it submits the inference, waits for completion, pushes a ReasonLink descriptor, and polls for the credit before reusing NVDLA's single output buffer.

This result depends on the deployment. Classifying all givens of a puzzle in one NVDLA submission would reduce $T_{\mathrm{nn}}$ substantially and allow one ReasonLink push per puzzle instead of one per cell; multiple output buffers would let NVDLA run ahead of the SAT-Accel tile. Under either change, conversion and submission costs that are currently hidden would become visible, and the comparison between CPU-generated and hardware-generated deltas could change. The replay diagnostic marks the opposite extreme: with inference removed, the faithful handoff is 99.2\% of instrumented runtime, but a full-image replay alone does not show where the two delta paths diverge. More expensive SAT search would instead amortize both neural and handoff work. Passing non-Sudoku functional checks does not constitute a performance study on harder workloads.

\subsection{Limits of the Current Evidence}
The evaluation uses one small structured workload family and 100 board puzzles per completed configuration. Sudoku is largely propagation dominated, so it is particularly suited to revealing loading and initialization overhead. Capacity limits, large learned-clause populations, sustained buffer backpressure, and long-running searches remain insufficiently explored. No power or energy measurements are available. Synthesized LUT comparisons are useful implementation context, but the current prototype does not establish an energy or area-efficiency advantage.

Measurement interference is another material limitation. A model-based subtraction can help diagnose the board's behavior, but raw timings and clean repeated runs must carry the final claims. The control path also retains per-cell CPU submissions and a DRAM round trip, so these results do not evaluate a fully autonomous direct-stream design.
